\chapter{델타에서 스냅샷을 재구성한다는 것}
\label{ch:orderbook}

이 장은 거래소가 매번 전체 호가창을 보내지 않고 변경분(delta)만 보내는 이유, 그리고
그 변경분을 누적해 전체 스냅샷(snapshot)을 다시 만들어내는 재구성기(reconstructor)가
왜 WS 연결을 직접 쥔 producer 안에 있어야만 하는지를 다룬다. 이 장에서 정리하는
\texttt{BookEvent}/\texttt{VenueOrderbookAdapter} 경계는 \Cref{ch:delivery}에서
다루는 순서 보장 문제(§4.4의 last-write-wins 가드)의 전제가 된다.

\todo{§4.0--§4.6을 바탕으로 재구성이 producer에 남아야 하는 이유(seq 갭 감지 시
reconnect 권한), Redis 라이터가 무상태인 이유(A안 채택 근거), venue 중립 어댑터
경계를 설명한다.}
